\section{Relation to earlier work}\label{sec:related}
\paragraph{Information and meaning.} Bar-Hillel and Carnap measured the semantic information of a statement by the possible worlds it excludes \citep{barhillel1953}; Floridi develops a philosophy of semantic information in which data must be well formed, meaningful and truthful \citep{floridi2011}. Both give a measure of a statement, not of a source, and no coding theorem. The ledger is a measure of a source, obtained from a joint law, and every entry has a coding meaning. \citet{kolchinsky2018} define semantic information as the information a system has about its environment that is causally necessary for its continued existence; the purpose $Y$ of \cref{def:purpose} plays the role of what matters to the receiver, but in our setting it is any variable the analyst names, and no physical theory of agency is used.

\paragraph{Semantic communication.} The engineering literature on semantic and goal-oriented communication is large \citep{bao2011,gunduz2023}. Formulations that ask the receiver to reproduce a function of the source are indirect rate--distortion problems, whose theory is due to \citet{dobrushin1962}, \citet{wolf1970} and \citet{witsenhausen1980} and which \citet{liu2021} apply to semantic sources with an intrinsic and an extrinsic state. \citet{goldreich2012} treat communication in which meaning arises from the goal of the exchange and the parties have no common language, and \citet{juba2011} show that ambiguity can be efficient when the parties do not share a prior. Our results on coding are consequences of Shannon's theorems and add nothing to these theories; what is added is the accounting that separates the three quantities every such formulation uses, and the leakage that tests a chosen meaning.

\paragraph{Meaning as sufficiency.} Defining the meaning of a message as the belief it induces is the construction of a minimal sufficient statistic \citep{fisher1922}, of the posterior as the canonical form of an experiment \citep{blackwell1953}, of the information bottleneck at the point of no loss \citep{tishby2000}, and of the causal state of computational mechanics \citep{shalizi2001}. Distributional clustering of words by the contexts they occur in \citep{pereira1993} and the distributional hypothesis behind it \citep{harris1954,firth1957} are the same idea applied to language. What \cref{thm:derived} and \cref{thm:leak} add is bookkeeping: the three-term split of the entropy into consequence, hedge and latitude, the monotonicity in the purpose, the leakage of a declared meaning, and the statement that no meaning is adequate for every purpose.

\paragraph{Ambiguity, synonymy and efficiency.} That languages trade the speaker's effort against the listener's is Zipf's principle of least effort \citep{zipf1949}; \citet{ferrer2003} put it in information-theoretic form with a cost that mixes the entropy of the signals and the information they carry. \citet{piantadosi2012} argue that ambiguity is efficient because context disambiguates, and \citet{gibson2019} review the efficiency view of language, and \citet{zaslavsky2018} give an information-bottleneck account of colour naming. The vocabulary law of \cref{cor:vocab} is an accounting statement inside this tradition: it says how much vocabulary a given content and freedom require, without a claim about what speakers optimise.

\paragraph{Semantic entropy of language models.} \citet{kuhn2023} and \citet{farquhar2024} cluster the answers of a language model by bidirectional entailment and take the entropy of the cluster distribution as an uncertainty score that flags likely confabulation. In the ledger's terms the clusters are a meaning map $\pi$ and their score is the content $H(\pi(X))$ of the model's output distribution; the ordinary sequence entropy is the expression $E$, and their difference is the latitude, the variation in wording that the entailment relation treats as irrelevant. The ledger gives the difference a name, and the leakage gives a test of whether the entailment clusters are adequate for a stated downstream use.

\paragraph{Clustering and distances between partitions.} $\Lat+\Dbt$ is the distance between partitions of \citet{rokhlin1967}, used in clustering as the variation of information \citep{meila2007} (\cref{prop:dist}). It is the sum of the two directions of conditional entropy, and the ledger keeps the directions apart.

\paragraph{Steganography and watermarking.} \cref{thm:channel} is the arithmetic-coding construction of provably secure steganography \citep{cachin1998,hopper2002}, used with language models by \citet{ziegler2019} and made efficient by minimum entropy coupling in \citet{dewitt2023}. Its capacity is the entropy of the cover given what must be preserved; the observation here is that ``what must be preserved'' can be the meaning, and that the capacity is then the latitude.

\paragraph{Codon usage.} The dependence of synonymous codon choice on base composition \citep{sueoka1962,muto1987}, on gene expression \citep{gouy1982,sharp1987} and its summary by an effective number of codons \citep{wright1990} or by Shannon entropy \citep{zeeberg2002} are classical. The genetic code has been treated as a channel before \citep{yockey1992}. We are not aware of an earlier split of the information a codon carries about its origin into levels of meaning by the chain rule, with the synonymous part identified as leakage; the numbers of \cref{sec:identity} are new in that sense, and nothing in them is surprising to a molecular biologist.
